% TOP_PROBLEM: {"id": 3143, "title": "57-regular Moore graph?", "category_id": 3, "category": {"id": 3, "name": "graph_theory", "display_name": "Graph Theory", "description": "Problems involving graphs, networks, and their properties.", "slug": "graph-theory", "order_index": 3, "created_at": "2026-07-31T15:26:25.671Z"}, "status": "open", "problem_number": "OPG-160", "set_id": 12}
% FIRST_POSTED: 2026-10-01
% ATTEMPT_STATUS: unresolved
% UnsolvedMath ID 3143; source catalogue code OPG-160
% !TeX program = lualatex
% Research attempt by GPT-6 Astra Ultra; no claim of a new complete solution.
\documentclass[11pt,a4paper]{article}
\usepackage[margin=25mm]{geometry}
\usepackage{fontspec}
\setmainfont{lmroman10-regular.otf}[BoldFont=lmroman10-bold.otf,
 ItalicFont=lmroman10-italic.otf,BoldItalicFont=lmroman10-bolditalic.otf]
\usepackage{amsmath,amssymb,mathtools}
\usepackage{unicode-math}
\setmathfont{latinmodern-math.otf}
\AtBeginDocument{\let\setminus\smallsetminus}
\usepackage{xurl,hyperref}
\hypersetup{colorlinks=true,urlcolor=blue}
\setcounter{secnumdepth}{0}
\setlength{\parindent}{0pt}
\setlength{\parskip}{5pt}
\setlength{\emergencystretch}{3em}
\begin{document}
\section*{57-regular Moore graph?}\label{3143}
{\small UnsolvedMath ID 3143 · OPG-160 · Graph Theory · OpenGarden}
\subsection{Definitions and mathematical statement}
Consider finite simple undirected graphs. The degree of a vertex is
its number of neighbours, and a graph is $57$-regular when every
vertex has degree $57$. In a connected graph the distance between
two vertices is the minimum number of edges in a joining path;
the diameter is the maximum of these distances. The girth is the
length of a shortest simple cycle.

The existence question asks whether there is a $57$-regular graph
$G$ of diameter two and girth five. Thus any two distinct vertices
must either be adjacent or have a common neighbour, while triangles
and quadrilaterals are forbidden. Equivalently, one seeks a strongly
regular graph with parameters
\[
 (v,k,\lambda,\mu)=(3250,57,0,1):
\]
there are $3250$ vertices, each has $57$ neighbours, adjacent
vertices have no common neighbour, and nonadjacent vertices have
exactly one common neighbour.

The order $3250$ follows by exploring distances from any vertex:
there are $1$ vertex at distance zero, $57$ at distance one, and
$57\cdot56$ distinct vertices at distance two. Girth five prevents
overlaps in that exploration, and diameter two exhausts the graph.
Hence the problem is a finite existence question at a prescribed
order, not a request for arbitrarily large graphs. Either an
explicit construction with the stated properties or a proof that
the parameter conditions are impossible would answer it.
\subsection{Short English statement}
Does a graph on 3250 vertices exist in which each vertex has 57 neighbours, adjacent pairs share none, and nonadjacent pairs share exactly one?
\subsection{Research attempt}
\paragraph{Scope and process.}
This is a GPT-6 Astra Ultra research attempt dated 1 October 2026, using
primary-source browsing and elementary mathematical checks. The canonical
definition above is retained. Results below are restricted results or
reductions, not a claim of a new solution to the entire catalogue question.
No formal proof assistant or external mathematical referee checked this text.


\paragraph{Literature and attack plan.}
The parameter set in [S2] is the missing degree-$57$ Moore graph.
A June 2026 preprint by Ishida [R1] studies restrictions on its
hypothetical automorphisms; its abstract claims that no involution can
occur. That claim is not used here, and excluding particular symmetries
would not exclude a graph with trivial automorphism group. Our plan is
instead to extract exact spectral restrictions and an equivalent system
of permutation-matrix equations, then check the reduction in degree three.

\paragraph{Spectral restrictions and a quantitative consequence.}
Let $A$ be the adjacency matrix of a putative graph and $J$ the all-one
matrix. Counting common neighbours gives
\[
 A^2=56I-A+J.
\]
The graph is connected, so its regular eigenvalue $57$ has multiplicity
one. On the orthogonal complement of the all-one vector, the remaining
eigenvalues satisfy $t^2+t-56=0$, and hence are $7$ and $-8$.
Writing their multiplicities as $f,g$, dimension and trace give
\[
 f+g=3249,\qquad 57+7f-8g=0,
 \qquad(f,g)=(1729,1520).
\]
Both multiplicities are positive integers; the spectral feasibility test
therefore supplies no contradiction.

There is nevertheless a useful restriction on independent sets. If $S$
is independent with size $a$, put $z=1_S-(a/3250)1$. Then
$z^TAz=-57a^2/3250$ and $z^Tz=a-a^2/3250$. The lower eigenvalue bound
$z^TAz\ge-8z^Tz$ yields $a\le400$. If $a=400$, equality forces
$Az=-8z$, so every vertex outside $S$ has precisely eight neighbours in
$S$. Thus even an extremal independent set would impose a regular
incidence structure rather than immediately contradicting the parameters.

\paragraph{An exact local-to-global matrix system.}
Fix a vertex $v$ and enumerate its neighbours $u_1,\ldots,u_{57}$.
For each $i$, let $B_i=N(u_i)\setminus\{v\}$. Girth five makes these
$57$ sets disjoint, independent, and of size $56$. They exhaust all
remaining vertices. A vertex $x\in B_i$ has exactly one neighbour in
each $B_j$ with $j\ne i$: the nonadjacent pair $x,u_j$ has exactly
one common neighbour, which must lie in $B_j$. The edges between
$B_i,B_j$ are consequently a perfect matching.

After labelling each block by $\{1,\ldots,56\}$, let $P_{ij}$ be its
matching matrix, with rows indexed by $B_i$ and columns by $B_j$.
Then $P_{ji}=P_{ij}^{T}$ and each $P_{ij}$ is a permutation matrix.
For distinct $i,j$, counting common neighbours of vertices in these two
blocks gives the entrywise equation
\[
 \boxed{\ \sum_{h\ne i,j}P_{ih}P_{hj}=J_{56}-P_{ij}.\ }
\]
There are $55$ summands. A matched pair has no common neighbour, while
every unmatched pair has one. These common neighbours can only occur
in the other blocks.

Conversely, suppose permutation matrices satisfying all these equations
exist. Add vertices $v,u_1,\ldots,u_{57}$ to the blocks, join $v$ to
all $u_i$, join $u_i$ to its entire block, and insert the specified
matchings. Every vertex has degree $57$. Two distinct vertices within
one block have their common neighbour $u_i$ and no other, since each
interblock matching is bijective. Pairs in different blocks have the
required common-neighbour count by the displayed equation. The remaining
pair types, involving $v$ or a vertex $u_i$, have counts zero when
adjacent and one otherwise, directly from the construction. This proves
that the matrix system is equivalent to the original existence problem.
It is a reduction, not a construction of its required solution.

\paragraph{Verification and first obstruction.}
For the degree-three analogue there are three blocks of size two. Set
$P_{12}=P_{13}=I_2$ and $P_{23}=J_2-I_2$, with reverse matrices their
transposes. Every equation holds. The resulting ten-vertex graph has
adjacency identity $A^2=2I-A+J$ and is a Petersen graph. The reproducible
script \texttt{scripts/check\_attempt\_batch02\_algebra\_2026\_10\_01.py}
checks every pair's degree and common-neighbour conditions for this
small construction. This validates the indices and diagonal exclusions
in the reduction; it supplies no evidence that the $56$-point system
is feasible.

\paragraph{Final status.}
\textbf{UNRESOLVED.} The spectral data are consistent, and the exact
permutation constraints retain all the difficulty. The first missing
step is either a family of $57$ mutually compatible matching systems
satisfying these equations, or a contradiction valid for every such
family without an unjustified symmetry assumption.

\subsection{Sources}
\begin{itemize}
\item[S1] ULAM.AI and the UnsolvedMath Contributors, supplied problems.json, record 3143 (OPG-160), OpenGarden collection. Catalogue identity and source transcription; CC BY 4.0.
\url{https://huggingface.co/datasets/ulamai/UnsolvedMath}.
\item[S2] Open Problem Garden, 57-regular Moore graph?. Original problem and discussion; the mathematical formulation below is expanded from the supplied source record.
\url{http://www.openproblemgarden.org/op/57_regular_moore_graph}.

\item[R1] Yawara Ishida, \emph{No involutions in the missing Moore graph},
2026, arXiv v2 dated 8 July 2026. Primary preprint abstract checked
1 October 2026; its automorphism claim is not used as a lemma.
\url{https://arxiv.org/abs/2606.29183}.

\end{itemize}
\end{document}
